\SourceLeaf{250}%
\noindent in lauter dreiseitige Pyramiden zerlegen, welche paarweise in beiden Körpern
symmetrisch sind.

II. Um jede dreiseitige Pyramide lässt sich eine Kugel construiren, deren
4 Halbmesser, mittelst durch je zwei gelegte Ebenen, dieselbe in je 4 gleichseitige
Pyramiden zerlegen. (Vorbehaltlich, dass auch negative Pyramiden dabei sein
können.)

III. Macht man diese Construction bei jedem Paar der ad I. erwähnten
dreiseitigen Pyramiden, so entstehen lauter gleichseitige Pyramiden, welche
wieder paarweise symmetrisch sind.

IV. Jede gleichseitige Pyramide lässt sich mittelst eines Perpendikels von
der Spitze auf die Grundfläche und durchgelegter Ebenen in \textls{drei} dreiseitige
Pyramiden zerfällen, deren Grundfläche ein gleichschenkliges Dreieck und
deren eine Kante, welche von der Spitze dieses gleichschenkligen Dreiecks
ausgeht, senkrecht auf der Grundfläche steht.

V. Zwei Pyramiden der letztbeschriebenen Art, die in beiden Polyedern
analog liegen, sind nicht nur symmetrisch, sondern auch congruent, weil sich
jede einzelne durch eine Ebene in zwei einander symmetrische Hälften spalten
lässt.

VI. Demnach sind beide Polyeder als die (algebraischen) Summen von lauter
congruenten Pyramiden zu betrachten, und haben also gleiches Volumen. Q.E.D.

Ich muss noch hinzufügen, dass der Beweis mir nicht ganz eigenthümlich ist.
Ich war im Nachdenken über Ihre Äusserung dahin gekommen, dass ich die
Theile III---V klar übersah, fand aber anfangs Schwierigkeit, jede dreiseitige
Pyramide auf gleichseitige zu reduciren; und hätte vielleicht die Sache gehen
lassen, wenn nicht in einem Gespräch mit Dr.\ \textsc{Stegmann} die Rede zufällig
darauf gekommen, wo dieser den glücklichen Einfall hatte, die Kugel um die
dreiseitige [Pyramide] zu Hülfe zu nehmen, wodurch dann der Schlüssel eigentlich
gefunden war.

So viel ich sehe, ist kein Fehlschluss in dem Obigen; Sie verbinden aber
mich sehr, wenn Sie den Beweis einer scharfen Prüfung unterwerfen und mir
gelegentlich Ihr Urtheil darüber schreiben. \mbox{.\ .\ .\ .\ .\ .\ .}\}
